\section{Generating opens and finite presentations}
\label{sec:finite-preliminaries}
Here ``generating'' refers to generation as an ideal or invertible
module. It does not assert generation as a unital algebra.

\begin{lemma}[The relative generating open]\label{lem:generating-open}
Let $B/S$ be a finite locally free commutative algebra over a
$\C$-scheme, and let $A\subset B$ be a subbundle. The following
conditions are equivalent: $B\otimes A\to B$ is surjective; on every
geometric fibre $A$ has nonzero image in every residue field of $B$;
and the open $A^\times$ of sections of $A$ which are units of $B$
is smooth and surjective over $S$. These conditions define an open
in the Grassmannian. The assertions have the same form for a
subbundle of an invertible $B$-module, using generating sections in
place of units.
\end{lemma}
\begin{proof}
Surjectivity is equivalent to its geometric-fibre version by
Nakayama's lemma. On a finite algebra over an algebraically closed
field, an ideal is the unit ideal exactly when its images in all
residue fields are nonzero. These evaluations cut out finitely many
proper hyperplanes in $A$. Their complement is nonempty over an
infinite field and consists precisely of units. On the total space
of the vector bundle $A$, invertibility is the open condition
$\det(b\mapsto ab)\ne0$. This open is smooth over $S$, and the
preceding fibre calculation proves surjectivity. Conversely a unit
section after a faithfully flat cover generates the unit ideal,
which descends. The Grassmannian open is the complement of the
support of the cokernel of $B\otimes A\to B$. Local trivializations
of an invertible module prove the last assertion.
\end{proof}

\begin{lemma}[Total-degree truncation]\label{lem:polarized-truncation}
Let $B$ be a rank-$d$ finite locally free algebra over a $\C$-algebra
$R$, and let $b_1,\ldots,b_q\in B$. The $R$-subalgebra they generate
is spanned as an $R$-module by the monomials of total degree at most
$d-1$. This assertion commutes with base change, with polynomial
reduction coefficients on a free chart.
\end{lemma}
\begin{proof}
Set $w=\sum_j t_jb_j$ in $B[t_1,\ldots,t_q]$ and apply
Cayley--Hamilton to multiplication by $w$. The coefficient of
$t^\alpha$, for $|\alpha|=d$, in $w^d$ is
$\binom d\alpha b^\alpha$. All other terms have $B$-monomial degree
less than $d$: coefficients of the characteristic polynomial are
scalars in $R[t]$. The multinomial coefficient is invertible in $R$.
This expresses each degree-$d$ monomial using lower degrees.
Multiplying these identities and inducting reduces every higher
monomial. The adjugate identity for Cayley--Hamilton makes the
coefficients polynomial in the multiplication table.
\end{proof}
